\subsection{Bias-variance-covariance-locality decomposition}
\label{subsec:biasvariance}
We now introduce our bias-variance-covariance-locality decomposition which extends the bias-variance decomposition \cite{kohavi1996bias} to WA.
In the rest of this theoretical section, $\ell$ is the Mean Squared Error for simplicity: yet, our results may be extended to other losses as in \cite{domingos2000unified}.
In this case, the expected error of a model with weights $\theta(l_S)$ \wrt the learning procedure $l_S$ was decomposed in \cite{kohavi1996bias} into:
\begin{equation} \tag{BV}\label{eq:b_v}
    \mathbb{E}_{l_S}\mathcal{E}_T(\theta(l_S)) = \mathbb{E}_{(x,y)\sim p_T}[\biasb^2(x, y)+\varb(x)],
\end{equation}
where $\biasb(x,y), \varb(x)$ are the bias and variance of the considered model \wrt a sample $(x,y)$, defined later in \Cref{eq:b_var_cov}.
To decompose WA's error, we leverage the similarity (already highlighted in \cite{izmailov2018}) between WA and functional ensembling (ENS) \cite{Lakshminarayanan2017,dietterich2000ensemble}, a more traditional way to combine a collection of weights.
More precisely, ENS averages the predictions, $f_{\text{ENS}} \triangleq f_{\text{ENS}}(\cdot, \{\theta_m\}_{m=1}^M) \triangleq 1/M \sum_{m=1}^{M} f(\cdot, \theta_m)$.
\Cref{lemma:wa_ensembling} establishes that $f_{\text{WA}}$ is a first-order approximation of $f_{\text{ENS}}$ when $\{\theta_m\}_{m=1}^M$ are close in the weight space.
\begin{lemma}[WA and ENS. Proof in \Cref{app:wa_loss}. Adapted from \cite{izmailov2018,Wortsman2022ModelSA}.] \label{lemma:wa_ensembling}
    Given $\{\theta_m\}_{m=1}^M$ with learning procedures $L_S^M\triangleq\{l_S^{(m)}\}_{m=1}^M$. Denoting $\Delta_{L_S^M}=\max_{m=1}^{M}\left\|\theta_m-\theta_{\text{WA}}\right\|_2$, $\forall (x,y) \in \X \times \Y$:%
    \begin{equation*}
        f_{\text{WA}}(x) = f_{\text{ENS}}(x) + O(\Delta^2_{L_S^M}) \text{~and~} \ell\left(f_{\text{WA}}(x), y\right) = \ell\left(f_{\text{ENS}}(x) , y\right) + O(\Delta_{L_S^M}^2).%
    \end{equation*}%
\end{lemma}%
This similarity is useful since \Cref{eq:b_v} was extended into a bias-variance-covariance decomposition for ENS in \cite{ueda1996generalization,brown2005between}.
We can then derive the following decomposition of WA's expected test error. To take into account the $M$ averaged weights, the expectation is over the joint distribution describing the $M$ identically distributed (i.d.) learning procedures $L_S^M\triangleq\{l_S^{(m)}\}_{m=1}^M$.
\begin{proposition}[Bias-variance-covariance-locality decomposition of the expected generalization error of WA in \ood. Proof in \Cref{app:proof_bvc}.]
    \label{prop:b_var_cov}
    Denoting $\bar{f}_S\left(x\right) = \mathbb{E}_{l_S} \left[f\left(x,\theta\left(l_S\right)\right)\right]$, under identically distributed learning procedures $L_S^M\triangleq\{l_S^{(m)}\}_{m=1}^M$, the expected generalization error on domain $T$ of $\theta_{\text{WA}}(L_S^M)\triangleq\frac{1}{M} \sum\nolimits_{m=1}^{M} \theta_m$ over the joint distribution of $L_S^M$ is:%
    \begin{equation}
        \begin{aligned}
             \mathbb{E}_{L_S^M}\mathcal{E}_T(\theta_{\text{WA}}(L_S^M)) &= \mathbb{E}_{(x,y)\sim p_T}\Big[\biasb^2(x, y)+\frac{1}{M} \varb(x)+\frac{M-1}{M} \covb(x)\Big] + O(\Deltab^2), \\
             \text{where~}\biasb(x,y)&=y-\bar{f}_S\left(x\right), \\
             \text{and~}\varb(x)&=\mathbb{E}_{l_S}\left[\left(f(x, \theta(l_S)) - \bar{f}_S\left(x\right)\right)^{2}\right], \\
             \text{and~}\covb(x)&= \mathbb{E}_{l_S,l_S'}\left[\left(f(x,\theta(l_S))-\bar{f}_S\left(x\right)\right)\left(f(x,\theta(l_S')))-\bar{f}_S\left(x\right)\right)\right], \\
             \text{and~}\Deltab^2&=\mathbb{E}_{L_S^M}\Delta_{L_S^M}^2 \text{~with~} \Delta_{L_S^M}=\max_{m=1}^{M}\left\|\theta_m-\theta_{\text{WA}}\right\|_2.%
        \end{aligned} \tag{BVCL}\label{eq:b_var_cov}%
    \end{equation}%
    $\covb$ is the prediction covariance between two member models whose weights are averaged.
    The locality term $\Deltab^2$ is the expected squared maximum distance between weights and their average.
\end{proposition}
\Cref{eq:b_var_cov} decomposes the \ood error of WA into four terms.
The bias is the same as that of each of its i.d. members.
WA's variance is split into the variance of each of its i.d. members divided by $M$ and a covariance term.
The last locality term constrains the weights to ensure the validity of our approximation.
In conclusion, combining $M$ models divides the variance by $M$ but introduces the covariance and locality terms which should be controlled along bias to guarantee low \ood error.